\section{Discussion and limitations}
\label{sec:discussion}
\paragraph{What the evidence supports.}
On the embedding leaderboards, a handful of well-chosen benchmarks recovers most of the variance in the other scores, and choosing them with the factor model is worth a quarter of the MSE at $k=3$ against random choice on A1 and (with random folds) A2; the advantage narrows as $k$ grows and depends on the criterion (the A-optimal criterion loses to random choice at large $k$ on A2 and HELM MMLU), the matrix and the imputer.
Small $k$ does not mean exact: a tight tolerance on RMSE, and recovery of the best models among the history, need many more benchmarks than the first factor suggests.

\paragraph{Limits.}
The error forecast is an optimistic reference rather than a calibrated one, and it needs a history several times larger than $N$; the factor model is less accurate than nonlinear imputers on MTEB, so point accuracy calls for pairing the selected set with one, and its intervals need the shift caveat of Figure~\ref{fig:f5}.
Our recommended criterion, the complement-error greedy, was added after a pilot and has no submodularity guarantee; on the sparse matrix B it is not better than $k$-medoids or forward regression, and its gain over random choice at $k=20$ there does not hold for held-out organisations; no stopping rule is validated under shift.
The embedding matrices are one family of leaderboards; the generative-LLM matrices have $N\le57$ benchmarks and mostly two or three factors, too few to carry a $k$-curve, so for them we claim only the comparisons of Section~\ref{sec:experiments}.
Code, pinned data extraction, seeds and result files are supplementary material; licences: Appendix~\ref{app:data}.
